\section{Introduction}

Melanoma accounts for a small share of skin-cancer diagnoses but the majority
of skin-cancer deaths \cite{siegel2023cancer}. Stage at detection largely
determines prognosis: the five-year relative survival is roughly 99\% for
localised disease but falls to about 35\% once the disease has spread to
distant sites \cite{siegel2023cancer}. Dermoscopy, which magnifies the lesion
surface and suppresses surface reflection, is the standard screening tool, yet
its interpretation is expertise-dependent and inter-observer agreement is
imperfect even among specialists.

Convolutional neural networks trained on dermoscopic and clinical images have
matched dermatologist-level accuracy in controlled reader studies
\cite{esteva2017dermatologist}, but their use in the clinic has not followed.
The recurring obstacle is opacity: a network that reports only a class
probability gives the clinician nothing to check against their own judgement,
and a prediction that cannot be interrogated is difficult to act on when the
stakes are a missed melanoma.

The ABCDE rule (Asymmetry, Border irregularity, Colour variation, Diameter,
Evolving lesion) is the clinical framework taught for melanoma assessment
\cite{abbasi2004early}, and dermatologists routinely use it to record and
communicate findings. A classifier whose output can be related back to these
familiar attributes is easier to audit than one that emits a probability alone,
because a reviewer can compare the model's stated evidence with what they see in
the image.

We present MelanomaNet, which augments a standard classifier with several
interpretability components that each address a different question a clinician
might ask:

\begin{itemize}
    \item \emph{Where did the model look?} GradCAM++ produces heatmaps over the
          image regions that most affected the prediction.
    \item \emph{Which clinical features are present?} An image-processing pipeline
          extracts the ABCDE criteria and reports quantitative scores.
    \item \emph{Which concepts drove the decision?} Concept activation vectors
          score the influence of clinically defined concepts on the prediction.
    \item \emph{How much should the output be trusted?} Monte Carlo dropout
          estimates predictive uncertainty and separates its epistemic and
          aleatoric parts.
\end{itemize}

The contributions of this paper are threefold. First, we integrate attention,
clinical-criterion, concept, and uncertainty explanations for a single
dermoscopy classifier and report them together on the same cases. Second, we
define alignment metrics that measure, against a random-attention control,
whether the model's attention falls on the lesion and its border rather than
the background. Third, we evaluate the attention maps for faithfulness using
deletion and insertion tests, so that the explanations are assessed rather than
merely displayed.
